In this section, we describe the experimental setup and the results obtained by comparing the introduced \varan to existing methods and techniques for fine-tuning SSL models on downstream tasks.

\begin{table*}[t]
    \caption{
The results of \varan are compared with other layer aggregation methods across various fine-tuning techniques. \varan outperforms alternative methods in most experiments and achieves comparable performance in a few cases. For clarity, the best results within the same backbone and fine-tuning method are highlighted in \textbf{bold}, while comparable performances are \underline{underlined}.
}
    \centering
    \begin{tabular}{l|l|cc|cc|ccc}
        \toprule
    \multicolumn{2}{c|}{\multirow{2}{*}{\textbf{Method}}} & \multicolumn{4}{c|}{\textbf{Speech Emotion Recognition}} & \multicolumn{3}{c}{\textbf{Automatic Speech Recognition}} \\
    
    \multicolumn{2}{c|}{} & \multicolumn{2}{c|}{\textbf{RAVDESS~\cite{livingstone2018ryerson}}} & \multicolumn{2}{c|}{\textbf{IEMOCAP~\cite{busso2008iemocap}}} & \textbf{GigaSpeech~\cite{DBLP:journals/corr/abs-2106-06909}} & \multicolumn{2}{c}{\textbf{LibriSpeech}~\cite{panayotov2015librispeech}} \\
  
    \multicolumn{2}{c|}{} & \shortstack{Acc. ($\uparrow$)} & \shortstack{Weighted \\ F1 ($\uparrow$)} & \shortstack{Acc. ($\uparrow$)} & \shortstack{Weighted \\ F1 ($\uparrow$)} & WER ($\downarrow$) & \shortstack{test-clean, \\ WER ($\downarrow$)} & \shortstack{test-other, \\ WER ($\downarrow$)} \\
    
    \midrule
    \multicolumn{9}{c}{\textbf{Data2Vec Base~\cite{baevski2022data2vec}}} \\
    \midrule
        \multirow{3}{*}{Fine-Tuning}
            & Last Layer   & 0.50 & 0.46 & 0.61 & 0.55 & \textbf{33.53} & \textbf{14.90} & \textbf{22.18} \\
            & Weighted Sum & 0.49 & 0.47 & \underline{0.65} & 0.63 & 79.20 & 61.69 & 75.56 \\
            & \varan       & \textbf{0.57} & \textbf{0.56} & \underline{0.65} & \textbf{0.65} & 42.93 & 25.32 & 31.97 \\ \midrule
        \multirow{3}{*}{LoRA} 
            & Last Layer   & 0.48 & 0.42 & \underline{0.71} & 0.69 & 37.92 & 20.76 & 26.81 \\
            & Weighted Sum & 0.52 & 0.49 & \underline{0.71} & 0.69  & 42.39 & 24.09 & 31.16 \\
            & \varan       & \textbf{0.60} & \textbf{0.60} & \underline{0.71}  & \textbf{0.70} & \textbf{36.04} & \textbf{18.71} & \textbf{24.67} \\
    \midrule
    \multicolumn{9}{c}{\textbf{Data2Vec Large~\cite{baevski2022data2vec}}} \\
    \midrule
        \multirow{3}{*}{Fine-Tuning} 
            & Last Layer   & 0.39 & 0.34 & \textbf{0.71} & \textbf{0.67} & 27.70 & \textbf{12.01} & \textbf{16.24} \\
            & Weighted Sum & 0.73 & 0.71 & 0.66  & 0.64 & 63.31 & 44.75 & 60.00 \\
            & \varan       & \textbf{0.79} & \textbf{0.79} & 0.66 & 0.63 & \textbf{27.53} & 12.90 & 17.89\\ \midrule
        \multirow{3}{*}{LoRA} 
            & Last Layer   & 0.40 & 0.36 & 0.69 & 0.65 & \textbf{27.12} & \textbf{13.38} & \textbf{17.40} \\
            & Weighted Sum & 0.66 & 0.65 & 0.70 & 0.69 & 34.39 & 18.71 & 25.01 \\
            & \varan       & \textbf{0.72} & \textbf{0.72} & \textbf{0.71} & \textbf{0.71} & 27.79 & 13.97 & 18.03 \\
    \midrule
    \multicolumn{9}{c}{\textbf{WavLM Base~\cite{chen2022wavlm}}} \\
    \midrule
        \multirow{3}{*}{Fine-Tuning} 
            & Last Layer   & 0.65 & 0.64 & \textbf{0.68} & \textbf{0.68}  & \textbf{50.73} & \textbf{31.92} & \textbf{44.19} \\
            & Weighted Sum & 0.64 & 0.64 & 0.65  & 0.62 & 67.35 & 49.33 & 65.46 \\
            & \varan       & \textbf{0.70} & \textbf{0.70} & 0.66 & 0.63 & 59.66 & 41.11 & 54.89 \\ \midrule
        \multirow{3}{*}{LoRA} 
            & Last Layer   & 0.48 & 0.48 & 0.70  & 0.68 & 48.70 & \textbf{29.20} & 39.94 \\
            & Weighted Sum & 0.54 & 0.54 & 0.66 & 0.65 & 52.56 & 33.60 & 44.00 \\
            & \varan       & \textbf{0.65} & \textbf{0.65} & \textbf{0.72} & \textbf{0.71} & \textbf{47.83} & 29.25 & \textbf{39.19} \\
    \midrule
    \multicolumn{9}{c}{\textbf{WavLM Large~\cite{chen2022wavlm}}} \\
    \midrule
        \multirow{3}{*}{Fine-Tuning}
            & Last Layer   & 0.46 & 0.42 & 0.69 & 0.68 & 31.62 & 18.59 & 26.34 \\
            & Weighted Sum & 0.73 & 0.71 & \underline{0.73} & \textbf{0.73} & 26.70 & \textbf{14.96} & \textbf{20.30} \\
            & \varan       & \textbf{0.75} & \textbf{0.76} & \underline{0.73} & 0.72 & \textbf{26.29} & 15.39 & 20.55 \\ \midrule
        \multirow{3}{*}{LoRA}
            & Last Layer   & 0.51 & 0.50 & 0.72 & 0.71 & 32.92 & 20.94 & 26.07 \\
            & Weighted Sum & \underline{0.60} & \underline{0.56} & 0.72 & 0.72 & 32.32 & 18.38 & 25.31 \\
            & \varan       & \underline{0.60} & \underline{0.56} & \textbf{0.73} & \textbf{0.73} & \textbf{27.39} & \textbf{15.55} & \textbf{21.17} \\
        \bottomrule
    \end{tabular}
    \label{table:main}
\end{table*}


\subsection{Setup}

To examine the \varan, we selected the Automatic Speech Recognition (ASR) and Speech Emotion Recognition (SER) tasks. These tasks require an understanding of acoustic features at both the semantic and phonetic levels, as well as speaker-dependent information. For the ASR task, we follow the experimental setup described in~\cite{zaiem2024less}. Specifically, we use the XS subset of the GigaSpeech dataset~\cite{DBLP:journals/corr/abs-2106-06909} for training. In addition to the test portion of the GigaSpeech dataset, we report results on the "test clean" and "test other" sets from the LibriSpeech dataset~\cite{panayotov2015librispeech}. For the SER task, we conduct experiments using the RAVDESS~\cite{livingstone2018ryerson} and IEMOCAP~\cite{busso2008iemocap} datasets. For the RAVDESS dataset, we use the original train-validation-test split and merge the neutral emotion with calm, resulting in a total of 7 training labels. For the IEMOCAP dataset, we perform 10-fold cross-validation by training on 9 speakers and testing on the remaining one.

As baselines for layer aggregation methods, we employ two approaches: (1) the straightforward \textit{last layer} method, where no aggregation occurs and features from the final layer are directly used; and (2) the \textit{weighted sum} method, where features are aggregated using learnable weights without conditioning on sample variation. Furthermore, we investigate different fine-tuning techniques. Initially, we keep the CNN feature extractor frozen and train the rest of the model's parameters on a downstream task, which can be effective but computationally expensive, we refer this setup as \textit{Fine-Tuning}. We also apply parameter-efficient fine-tuning using LoRA~\cite{hu2021lora}, a method that is both fast and stable for training large SSL models on downstream tasks, particularly under limited computational resources. More specifically, we follow~\cite{zaiem2024less} by replacing the feed-forward layers that come after the self-attention mechanism with LoRA layers.

For the backbones, we use both WavLM~\cite{chen2022wavlm} and Data2Vec~\cite{baevski2022data2vec} in base and large configurations. For both ASR and SER tasks, we used 2 fully-connected layers with a hidden size of 256 as a probing head. ASR decoding was conducted at the character level using beam search (beam size=5, beam threshold=20), with no language model integration.

Since \varan’s training objective derives from a variational upper bound (see Section~\ref{sec:method}), one essential consideration is selecting an appropriate prior distribution. We adopt a discretized reversed $\chi^2$ as the prior, with degrees of freedom treated as hyperparameters during search.

For each fine-tuning, we use grid search to find the optimal hyperparameters on the validation sets of the corresponding datasets. See Table \ref{tab:hp} for hyperparameter grid. For the IEMOCAP dataset, "Session 5" is used. Full grid of the hyperparameters can be found in Table \ref{tab:hp}.

\subsection{Results}
\label{sec:4_2}

Experimental results are summarized in Table~\ref{table:main}. On speech emotion recognition, \varan consistently outperforms other layer aggregation methods across configurations. The sole exception occurs in the LoRA fine-tuning setup with WavLM Large, where \varan achieves performance comparable to weighted-sum aggregation.

For the IEMOCAP dataset, \varan surpasses conventional methods less frequently in \textit{Fine-Tuning} setup. However, the strongest results on this benchmark are observed in the LoRA fine-tuning regime, where \varan delivers substantial improvements over all baselines, achieving a relative accuracy gain of 2.87\% compared to weighted-sum aggregation. 

According to evaluations on the GigaSpeech ASR corpus and LibiSpeech test-other subsets, \varan outperforms other methods in the LoRA fine-tuning setup across all models except Data2Vec Large. The relative improvement in WER on the GigaSpeech test set is 7.7\% compared to the last-layer baseline. In the \textit{Fine-Tuning} setup, performance gains are less stable: the last layer predominantly outperforms other methods.

This phenomenon can be explained by two observations. First, ~\cite{toyota23} demonstrated that models trained to predict discrete units learn phonetic and word information concentrated in higher layers. Second, during LoRA fine-tuning, models retain more pretraining knowledge~\cite{zaiem2024less}, making layer aggregation in a data-dependent setup beneficial. Conversely, in \textit{Fine-Tuning}, knowledge distribution may collapse to the last layer, diminishing the utility of multi-layer aggregation.

In general, \varan consistently outperforms other methods in most scenarios, particularly in the LoRA setup and the speech emotion recognition task, where low-level features and layer weighting play a crucial role in final performance.

\subsection{Analysis}
\label{sec:analysis}

To investigate the differences between datasets and models we plot the PMF function of prior distributions $q(l)$ found during hyperparameter search for \varan models. The results are presented in Figure \ref{fig:prior}. The first observation is that both Data2Vec \cite{baevski2022data2vec} base and large models benefit from the weights skewed towards the last layers in contrast to WavLM \cite{chen2022wavlm}, which indicates that features needed for the tasks are grouped there. The second observation is that ASR benefits most from the last hidden state.
